%!TEX root = ../main.tex
\section{Experiments and Results}
\label{sec:results}

We evaluate eight vision-language models across 250 scientific figures and over 14 experimental conditions, yielding more than 11,000 individual evaluation instances. The central finding is that model rankings under perception, reasoning, and behavioural evaluation are not the same ranking. GPT-5.2 produces the highest-quality descriptions in our benchmark (MQM 91.6) yet fabricates answers for elements it cannot see 94\% of the time, with nothing in its output to distinguish the fabrication from genuine analysis.

\subsection{Models}
\label{sec:models}

We select eight models spanning commercial and open-weight families, dense and mixture-of-experts (MoE) architectures, and a range of effective parameter counts.
\textbf{Commercial models}: GPT-5.2 (Azure, closed-weight), Gemini 3.1 Pro (closed-weight), and Phi-4-Multimodal (small, closed-weight).
\textbf{Open-weight MoE models}: Llama 4 Maverick (17B params, 128 experts), Qwen3-VL-235B-A22B (235B total, 22B active), and Qwen3-VL-30B-A3B (30B total, 3B active).
\textbf{Open-weight dense models}: Qwen3-VL-8B and Gemma3-27B-IT.
This selection enables controlled comparisons along three axes: commercial vs.\ open-weight deployment, MoE vs.\ dense architecture at similar active parameter counts (Qwen-235B at 22B active vs.\ Gemma at 27B dense), and scaling within a single family (the three Qwen variants from 3B to 22B active parameters). All models received identical prompts and were evaluated at temperature~0 with GPT-4o as the automated judge.

\subsection{Perception}
\label{sec:perception-results}

\paragraph{Baseline description quality.}
GPT-5.2 leads with a baseline MQM of 91.6 [90.4, 92.8], followed by Gemini at 90.2 [88.9, 91.4] (Table~\ref{tab:mqm-results}). The 1.4-point gap is statistically significant ($p < 0.01$) but practically small (Cliff's $\delta = 0.09$, negligible). Below these two leaders, a clear tier structure emerges. Llama~4, Qwen-235B, and Qwen-8B cluster between 78.9 and 81.4 with overlapping confidence intervals and negligible pairwise effect sizes. Qwen-30B (74.4) and Gemma (69.1) form a third tier. Phi-4 trails at 62.2 [59.3, 64.9], nearly 30 points behind the leaders, with completeness penalties (8.8 points) driving the gap far more than accuracy penalties do for stronger models (GPT-5.2 loses only 0.7 points to completeness).

\paragraph{Transform robustness.}
Rotation is the most damaging transform, causing an average MQM drop of 19.4 points across models, while noise has negligible impact (Table~\ref{tab:mqm-results}). Low contrast reduces scores by a moderate 4--7 points for top-tier models. The most revealing condition is in-paper-blur, where the figure is embedded in a PDF page and the surrounding content is blurred. This condition collapses all models below 31 MQM points. Gemini drops from 90.2 to 8.4. GPT-5.2 drops from 91.6 to 13.3. Both attempt to describe the blurred page context rather than focusing on the embedded figure, a failure with practical implications for document-level figure understanding pipelines.

\paragraph{Selective blur.}
Selectively blurring unrecoverable elements (admittance blur) reduces MQM by roughly 8--10 points for top models, while blurring inferable elements (inductance blur) produces smaller drops of 1--3 points (Table~\ref{tab:mqm-results}, AdmB and IndB columns). This difference validates the admittance-inductance distinction at the perception level. GPT-5.2 scores 88.9 under inductance blur but 81.7 under admittance blur, a 7.2-point gap that reflects the difference between elements a model can reason about and elements that are simply gone.

\paragraph{Caption bias on description quality.}
Poisoned captions have no measurable effect on MQM for resistant models. GPT-5.2 and Gemini both score 91.3 under modified captions, matching their baselines, because they ignore the false claims entirely. Phi-4, already caption-dependent, scores 61.7 under the same condition, consistent with its near-zero resistance ($R = 0.05$).

\subsection{Reasoning}
\label{sec:reasoning-results}

Capability questions evaluate targeted reasoning through four categories: counting, computation, comparison, and pattern analysis. Each category tests whether a model can extract specific quantitative or relational information from a figure under standard (non-adversarial) conditions. Full results appear in Table~\ref{tab:capability}.\footnote{Capability evaluation is currently being finalized across all eight models; we report preliminary results in the appendix and will include the complete table in the final version.}

\subsection{Behaviour}
\label{sec:behaviour-results}

Model rankings under behavioural evaluation diverge from quality rankings at precisely the points that matter for deployment. GPT-5.2 ranks first on description quality but falls to fourth on active admittance, while Gemini leads on nearly every behavioural dimension despite placing second on quality (Table~\ref{tab:behavioral}).

\paragraph{Resistance to hallucination probes.}
Gemini achieves the highest overall resistance at 0.91 [0.89, 0.93], followed by GPT-5.2 at 0.81 [0.79, 0.84]. Phi-4 anchors the bottom at 0.21 [0.18, 0.24]. The gap between Gemini and GPT-5.2 is statistically significant ($p < 0.001$, $n = 750$ probe instances), confirming that resistance is not simply a correlate of overall model capability.

Probe type difficulty follows a clear gradient. Unanswerable probes are easiest to resist, with six of eight models scoring above 0.88, because the question format permits straightforward refusal. Inexist probes, grounded in presupposition embedding \citep{loftus1975leading}, are hardest, with scores ranging from 0.04 (Phi-4) to 0.88 (Gemini). This 22$\times$ gap reflects the power of definite articles to bypass model vigilance. A model that can refuse an explicitly unanswerable question still accepts the existence of chart elements that were never there. Contra probes fall between these extremes, with anchoring bias \citep{tversky1974judgment} proving moderately effective at suppressing model pushback.

\paragraph{Caption bias resistance.}
Models fall along a caption dependency spectrum from near-total visual independence to near-total textual dependency (Table~\ref{tab:behavioral}, Cap.\ Bias~R). Gemini and GPT-5.2 both achieve 0.89 [0.85, 0.93], with overlapping CIs and a non-significant difference ($p = 0.44$). Llama~4 provides moderate resistance at 0.74. Phi-4 exhibits near-zero resistance at 0.05 [0.02, 0.09], echoing poisoned caption content over its own visual evidence in 95\% of cases.

Caption dependency tracks active parameter count within the Qwen family: 0.43 (8B) $\rightarrow$ 0.30 (30B, 3B active) $\rightarrow$ 0.54 (235B, 22B active). The non-monotonic pattern suggests that active parameter count, rather than total parameter count, predicts caption independence for MoE architectures.

\paragraph{Active admittance and inductance.}
Gemini is the only model that consistently acknowledges visual limitations, admitting uncertainty in 90\% of cases [82\%, 98\%] when asked about selectively blurred elements (Table~\ref{tab:behavioral}). No other model exceeds 22\%. GPT-5.2, the top-ranked model on description quality, admits visual limitations only 6\% of the time while fabricating answers 98\% of the time. This ``confident fabricator'' profile, high descriptive fluency coexisting with near-total unwillingness to acknowledge uncertainty, represents a concrete safety risk for scientific deployment.

The inductance dimension validates the A-R-I framework empirically. When models fabricate answers for \emph{inferable} elements (inductance probes), correctness ranges from 21\% to 80\% across models (Table~\ref{tab:behavioral}). When they fabricate for \emph{unrecoverable} elements (admittance probes), correctness drops to 0--14\%, consistent with lucky guessing. GPT-5.2 achieves 74\% inductance correctness ($n = 47$ fabrications) versus 14\% admittance correctness ($n = 49$), a 60-percentage-point gap that confirms models perform genuine contextual reasoning when context permits rather than generating plausible-sounding noise.

\paragraph{Passive admittance.}
Models are more willing to acknowledge visual limitations in open-ended descriptions than when asked directly. GPT-5.2 mentions blur or illegibility for 26\% of admittance-blurred elements in its descriptions but admits uncertainty for only 6\% when asked a targeted question (Table~\ref{tab:behavioral}). Gemini shows the same pattern at a higher baseline (74\% passive vs.\ 90\% active). This discrepancy suggests that targeted questions create social pressure to provide a definitive answer, consistent with Grice's cooperative principle \citep{grice1975logic}, while open-ended descriptions provide more natural opportunities for hedging.

\subsection{Ablation: Probe Designer Independence}
\label{sec:ablation}

A potential concern with LLM-generated probes is circularity. If the same model family designs the probes and evaluates the responses, results may reflect shared biases rather than genuine model behaviour. We compare probes designed by GPT-4o (our default) against probes designed by Mistral Large, evaluating both on the same 50-figure subset with GPT-5.2 as the target model. Caption bias resistance is nearly identical: 0.89 (GPT-4o probes) vs.\ 0.89 (Mistral probes), with a negligible effect size ($\delta = -0.01$, $p = 0.49$). Hallucination resistance shows a small difference (0.80 vs.\ 0.86, $\delta = -0.16$) whose 95\% CI includes zero. These results indicate that the behavioural patterns we observe are properties of the evaluated models, not artifacts of the probe designer.
